\hypertarget{appendix-1a50e72435c2}{}
\section*{Definite integrals and the fundamental theorem}
\chapter{Definite integrals and the fundamental theorem}\label{appendix-1a50e72435c2-definite-integrals-and-the-fundamental-theorem}

\textbf{Curriculum depth:} Developed · \textbf{First seen:} Phase 9 · \textbf{Latest development:} Phase 9

\section{Definition and mechanism}\label{appendix-1a50e72435c2-definition-and-mechanism}

A definite integral accumulates a quantity over an interval, with sign and units determined by the integrand and integration variable. Riemann sums approximate this accumulation by weighted samples; under suitable regularity conditions their limit defines the integral. The fundamental theorem connects accumulation and differentiation: an antiderivative evaluates the net integral, and the derivative of an accumulated integral recovers its integrand when continuity conditions hold. This link explains why integrating velocity gives displacement and why a variable upper-limit integral solves a simple differential equation. Definite integral, indefinite antiderivative, and geometric area are related but not interchangeable ideas. The integration variable is a dummy variable, so name it distinctly from the endpoint. Integrals also form the target of Monte Carlo estimators, where measure and sampling density must match.

\section{Prerequisites}\label{appendix-1a50e72435c2-prerequisites}

\begin{itemize}
\tightlist
\item
  \hyperref[appendix-7bc181935c42]{Derivative as a local linear map}
\end{itemize}

\section{Used by}\label{appendix-1a50e72435c2-used-by}

\begin{itemize}
\tightlist
\item
  \hyperref[appendix-80521eae6b71]{Numerical quadrature and error}
\item
  \hyperref[appendix-e8ab9836882e]{ODE initial-value problems and stiffness}
\end{itemize}

\section{Full teaching explanation}\label{appendix-1a50e72435c2-full-teaching-explanation}

\begin{itemize}
\tightlist
\item
  \href{Continuous\%20Notes/Computing\%20Foundations\%20-\%20From\%20Bits\%20to\%20C++\%20and\%20Systems\%20APIs\#Integration\%20and\%20quadrature}{Phase 9 teaching section}
\item
  \href{../01-Foundations-I/main.pdf\#nameddest=appendix-adf48f5dc5c1}{Visual Recall/09 - Calculus and Optimization - Visual Recall.md (01-Foundations-I, appendix-adf48f5dc5c1)} · visual recall
\end{itemize}

\section{Source coverage}\label{appendix-1a50e72435c2-source-coverage}

Supplied lesson parts: 832 (T03 · M051), 833 (T03 · M051), 834 (T03 · M051).
